%% 自動生成（tools/mknext.py）。各章の道しるべ「この章の先に」から、後ろの巻の章を集めたもの
\begin{longtblr}{colspec={Q[l,m,wd=58mm] X[l,m]}, width=\linewidth, rowhead=1, row{1}={bg=bkVol!14, font=\sffamily\bfseries\footnotesize},
  hlines={0.4pt,bkLine}, rowsep=1.5pt, cells={font=\footnotesize}}
この\R{巻}{かん}の\R{章}{しょう} & この\R{道具}{どう|ぐ}を\R{使}{つか}う\R{章}{しょう} \\
{\bknextpill{2-2}\ マイナスの\R{数}{かず}と\R{数直線}{すう|ちょく|せん}} & {\bknextpill{3-2}\ \R{場所}{ば|しょ}を\R{数}{かず}で\R{表}{あらわ}す：\R{座標}{ざ|ひょう}\\ \R{第}{だい}4\R{巻}{かん}から\R{先}{さき}：\bknextpill{4-10}\,\bknextpill{5-5}} \\
{\bknextpill{2-3}\ ＋ － × ÷ は、\R{何}{なに}を\R{表}{あらわ}しているの？} & {\bknextpill{3-5}\ グラフの\R{傾}{かたむ}きは「\R{変化}{へん|か}の\R{速}{はや}さ」\\ \R{第}{だい}4\R{巻}{かん}から\R{先}{さき}：\bknextpill{5-1}\,\bknextpill{5-9}} \\
{\bknextpill{2-5}\ \R{比}{ひ}・\R{割合}{わり|あい}・\R{百分率}{ひゃく|ぶん|りつ}} & {\R{第}{だい}4\R{巻}{かん}から\R{先}{さき}：\bknextpill{5-6}} \\
{\bknextpill{2-6}\ はかった\R{数}{かず}には「ずれ」がある} & {\bknextpill{3-14}\ ばらつきを1つの\R{数}{かず}で：\R{標準偏差}{ひょう|じゅん|へん|さ}\\ \R{第}{だい}4\R{巻}{かん}から\R{先}{さき}：\bknextpill{5-8}} \\
{\bknextpill{2-7}\ \R{何回}{なん|かい}もかける：\R{累乗}{るい|じょう}と\R{指数}{し|すう}} & {\R{第}{だい}4\R{巻}{かん}から\R{先}{さき}：\bknextpill{5-14}} \\
{\bknextpill{2-8}\ 2\R{回}{かい}かけると□になる\R{数}{かず}：\R{平方根}{へい|ほう|こん}} & {\bknextpill{3-10}\ ボールはいつその\R{高}{たか}さに\R{届}{とど}く？：2\R{次方程式}{じ|ほう|てい|しき}\\ \R{第}{だい}4\R{巻}{かん}から\R{先}{さき}：\bknextpill{4-4}\,\bknextpill{6-6}} \\
{\bknextpill{2-9}\ \R{分数}{ぶん|すう}で\R{書}{か}けない\R{数}{かず}がある？} & {\R{第}{だい}4\R{巻}{かん}から\R{先}{さき}：\bknextpill{4-2}\,\bknextpill{4-4}\,\bknextpill{5-14}} \\
{\bknextpill{2-10}\ \R{単位}{たん|い}の\R{世界}{せ|かい}ルール：SI\R{単位}{たん|い}} & {\R{第}{だい}4\R{巻}{かん}から\R{先}{さき}：\bknextpill{5-1}\,\bknextpill{5-5}} \\
{\bknextpill{2-11}\ なぜ\R{文字}{も|じ}を\R{使}{つか}うの？} & {\bknextpill{3-3}\ \R{関数}{かん|すう}は「\R{入}{い}れると\R{出}{で}てくる」ルール} \\
{\bknextpill{2-12}\ \R{文字式}{も|じ|しき}のやくそく} & {\bknextpill{3-8}\ かっこをはずす：\R{展開}{てん|かい}と\R{平方完成}{へい|ほう|かん|せい}\\ \R{第}{だい}4\R{巻}{かん}から\R{先}{さき}：\bknextpill{5-2}\,\bknextpill{5-4}} \\
{\bknextpill{2-13}\ \R{公式}{こう|しき}に\R{数}{かず}を\R{入}{い}れてみよう：\R{代入}{だい|にゅう}} & {\bknextpill{3-3}\ \R{関数}{かん|すう}は「\R{入}{い}れると\R{出}{で}てくる」ルール\\ \R{第}{だい}4\R{巻}{かん}から\R{先}{さき}：\bknextpill{5-2}\,\bknextpill{5-3}} \\
{\bknextpill{2-14}\ \R{式}{しき}を\R{変形}{へん|けい}して、ほしい\R{量}{りょう}をとり\R{出}{だ}そう} & {\bknextpill{3-6}\ 2\R{本}{ほん}の\R{直線}{ちょく|せん}が\R{交}{まじ}わるところ：\R{連立方程式}{れん|りつ|ほう|てい|しき}\\ \R{第}{だい}4\R{巻}{かん}から\R{先}{さき}：\bknextpill{5-0}} \\
{\bknextpill{2-15}\ \R{比例}{ひ|れい}と\R{反比例}{はん|ぴ|れい}} & {\bknextpill{3-3}\ \R{関数}{かん|すう}は「\R{入}{い}れると\R{出}{で}てくる」ルール\\ \bknextpill{3-4}\ グラフをかこう・\R{読}{よ}もう\\ \bknextpill{3-7}\ \R{反比例}{はん|ぴ|れい}のグラフ\\ \R{第}{だい}4\R{巻}{かん}から\R{先}{さき}：\bknextpill{5-12}\,\bknextpill{6-6}} \\
{\bknextpill{2-16}\ \R{何秒後}{なん|びょう|ご}に\R{追}{お}いつく？1\R{次方程式}{じ|ほう|てい|しき}} & {\bknextpill{3-4}\ グラフをかこう・\R{読}{よ}もう\\ \bknextpill{3-6}\ 2\R{本}{ほん}の\R{直線}{ちょく|せん}が\R{交}{まじ}わるところ：\R{連立方程式}{れん|りつ|ほう|てい|しき}\\ \bknextpill{3-10}\ ボールはいつその\R{高}{たか}さに\R{届}{とど}く？：2\R{次方程式}{じ|ほう|てい|しき}\\ \R{第}{だい}4\R{巻}{かん}から\R{先}{さき}：\bknextpill{6-2}} \\
{\bknextpill{2-17}\ コンピュータの\R{数}{かず}：2\R{進数}{しん|すう}と16\R{進数}{しん|すう}} & {\bknextpill{3-11}\ \R{指数関数}{し|すう|かん|すう}：いまある\R{量}{りょう}に\R{比例}{ひ|れい}してふえる・へる} \\
{\bknextpill{2-18}\ \R{表計算}{ひょう|けい|さん}で\R{計算}{けい|さん}しよう} & {\R{第}{だい}4\R{巻}{かん}から\R{先}{さき}：\bknextpill{5-8}\,\bknextpill{6-2}} \\
{\bknextpill{2-19}\ プログラムで\R{計算}{けい|さん}しよう} & {\R{第}{だい}4\R{巻}{かん}から\R{先}{さき}：\bknextpill{5-2}\,\bknextpill{6-2}} \\
\end{longtblr}
